\documentclass[11pt,a4paper]{article}
\usepackage[T1]{fontenc}
\usepackage[utf8]{inputenc}
\usepackage[ngerman]{babel}
\usepackage[margin=2.5cm]{geometry}
\usepackage[hidelinks]{hyperref}
\setlength{\parindent}{0pt}
\begin{document}
\begin{center}
  {\Large\textbf{Webtechnologien -- Übungsblatt 1}}
\end{center}
\section*{Aufgabe 3: HTML Page 1}
\textbf{Aufgabe:} Die vorgegebene Webseite ausschließlich mit HTML
nachbauen und über Apache testen.

\medskip
\textbf{Ergebnis:} Die Datei \texttt{index.html} enthält:
\begin{itemize}
  \item HTML-Grundgerüst mit Seitentitel, UTF-8 und Viewport-Meta-Tag.
  \item Überschrift, Absätze und zwei horizontale Linien.
  \item AAU-Logo direkt von der AAU-Webseite, 150\,$\times$\,55 Pixel.
  \item Drei Listenpunkte mit jeweils zwei nummerierten Unterpunkten.
  \item AAU-Link, der in einem neuen Tab geöffnet wird.
\end{itemize}
\textbf{Test:} Die Datei liegt unter
\path{C:/xampp/htdocs/webtech-aufgabe3/index.html} und ist über
\url{http://localhost/webtech-aufgabe3/} erreichbar.
Apache liefert die Seite mit Status \texttt{200 OK} aus.
\medskip
{\small \textbf{Hilfsmittel:} Codex für Erstellung, Prüfung und Dokumentation.}
\end{document}
